\documentclass[11pt]{article}
\usepackage[a4paper,margin=25mm]{geometry}
\usepackage{amsmath,amssymb,amsthm}
\usepackage[hidelinks]{hyperref}
\newtheorem{theorem}{Theorem}
\setlength{\emergencystretch}{1em}
\begin{document}
\title{A singleton cyclic-sieving instance for Conjecture 426}
\author{Chengzhe Gao}\date{}\maketitle
\begin{theorem}
There is a cyclic-sieving triple with a prime-order rotation group and
a polynomial that is equal to $1$ at every root of unity of that order.
\end{theorem}
\begin{proof}
Take two cyclically arranged positions and let $C=C_2$ act by their
rotation. This is a group of order $2$, which is prime. Let $X$ consist
of the single binary word $00$, and put $X(q)=1$. Rotating either zero
or one step fixes this word. Thus, for every $g\in C$,
\[
 |X^g|=1=X(\omega(g)), \qquad \omega(1_C)=1,\quad
 \omega(\text{one-step rotation})=-1.
\]
The two values of $\omega$ are precisely the complex second roots of
unity, and $X(q)$ has nonnegative integer coefficients. These are all
the required cyclic-sieving equalities. Indeed, the polynomial is
equal to $1$ at every complex number.
\end{proof}

The group rotates the two positions faithfully, while its action on
$X$ is trivial. The usual definition of a cyclic-sieving triple does
not require a faithful action on $X$ \cite{RSW}. This distinction is
essential here and is not concealed by the construction.

\paragraph{Scope of the statement.}
Both language versions of the conjecture ask for existence and impose
no condition $|X|>1$. In fact that additional condition would be
impossible: $1$ is one of the $p$-th roots of unity and the identity
element fixes every member of $X$, so the two stipulated equalities
give $|X|=X(1)=1$. Thus the parenthetical description of a separation
from ``nontrivial counting'' cannot mean a cardinality greater than
one. We prove the explicit existence assertion as written and also
explain this unavoidable limitation.

\paragraph{Formal verification.}
The Lean project defines the two-element cyclic group by Boolean
exclusive-or and checks its full group laws, enumeration, generator,
and prime order. Words are actual functions from the two positions to
the binary alphabet. The action is actual permutation of positions;
closure of $X$ and both action laws are proved. A Boolean equality
test is proved equivalent to function equality, so
\texttt{fixedCount} counts the genuine fixed words of the complete,
duplicate-free list $X$.

The polynomial is the coefficient list $[1]$, evaluated by Horner's
rule. The theorem \texttt{polynomial\_\allowbreak{}one\_\allowbreak{}at\_\allowbreak{}every\_\allowbreak{}point} holds
over every type supplied with the elementary evaluation operations
and the laws $q0=0$, $q+0=q$, and the natural-number image of $1$.
In particular it holds over the complex numbers, with no restriction
on the evaluation point. Hence this is not an evaluation at a finite
sample of roots. The concrete root map into the integers checks
that $1,-1$ multiply according to $C_2$ and are distinct; the same
values embed into the complex numbers. The universal evaluation
theorem proves all cyclic-sieving equalities under that embedding.
The file also proves that the identity equality forces cardinality
one. No custom axioms, admitted proofs, or external computation
oracles are used.

\begin{thebibliography}{1}
\bibitem{RSW} V. Reiner, D. Stanton and D. White,
\emph{The cyclic sieving phenomenon}, J. Combin. Theory Ser. A
108 (2004), 17--50, Definition--Proposition in Section 1.
\href{https://doi.org/10.1016/j.jcta.2004.04.009}{doi:10.1016/j.jcta.2004.04.009}.
\end{thebibliography}
\end{document}
